\lyseis{2.5}{Πολλαπλασιασμός κλασμάτων}
\sol{1} α) $\frac{10}{21}$ \quad β) $\frac{24}{36}=\frac{2}{3}$ \quad γ) $\frac{60}{150}=\frac{2}{5}$ \quad δ) $\frac{140}{210}=\frac{2}{3}$ \quad ε) $\frac{12}{72}=\frac{1}{6}$ \quad στ) $\frac{63}{168}=\frac{3}{8}$.
\sol{2} α) $\frac{8}{5}$ \quad β) $\frac{3\cdot 14}{7}=6$ \quad γ) $\frac{30}{9}=\frac{10}{3}$ \quad δ) $\frac{36}{8}=\frac{9}{2}$ \quad ε) $\frac{5\cdot 18}{6}=15$ \quad στ) $\frac{700}{10}=70$.
\sol{3} α) $\frac{1}{4}$ \quad β) $\frac{1}{15}$ \quad γ) $\frac{2}{3}\cdot\frac{3}{4}=\frac{1}{2}$ \quad δ) $\frac{3}{4}\cdot\frac{4}{5}=\frac{3}{5}$.
\sol{4} α) $\frac{7}{3}\cdot 3=7$ \quad β) $\frac{3}{2}\cdot\frac{8}{3}=4$ \quad γ) $\frac{13}{4}\cdot 2=\frac{13}{2}=6\frac{1}{2}$ \quad δ) $\frac{5}{2}\cdot\frac{8}{5}=4$.
\sol{5} $\frac{5}{7}$, $\frac{1}{3}$, $8$, $1$, $\frac{4}{11}$, και για το $2\frac{1}{2}=\frac{5}{2}$ ο αντίστροφος είναι $\frac{2}{5}$.
\sol{6} α) $\frac{4}{3}$ \quad β) $\frac{9}{5}$ \quad γ) $\frac{1}{6}$ \quad δ) $\frac{7}{2}$.
\sol{7} α) $\frac{14\cdot 25}{15\cdot 28}=\frac{350}{420}=\frac{5}{6}$ \quad β) $\frac{378}{945}=\frac{2}{5}$ \quad γ) $\frac{120}{144}=\frac{5}{6}$ \quad δ) $\frac{140}{252}=\frac{5}{9}$.
\sol{8} α) $\frac{2}{3}+\frac{3}{8}=\frac{16}{24}+\frac{9}{24}=\frac{25}{24}=1\frac{1}{24}$ \quad β) $\frac{7}{6}\cdot\frac{3}{4}=\frac{21}{24}=\frac{7}{8}$ \quad γ) $\frac{1}{2}\cdot\left(\frac{9}{12}-\frac{2}{12}\right)=\frac{1}{2}\cdot\frac{7}{12}=\frac{7}{24}$.
\sol{9} α) Σωστό. β) Λάθος: $\frac{3}{5}\cdot 4=\frac{12}{5}$. γ) Λάθος: $\frac{1}{2}\cdot\frac{1}{3}=\frac{1}{6}$. δ) Λάθος: το 0 δεν έχει αντίστροφο. ε) Σωστό. στ) Λάθος: $2\frac{1}{2}\cdot 2=\frac{5}{2}\cdot 2=5$.
\sol{10} $\frac{3}{4}\cdot 2\frac{1}{2}=\frac{3}{4}\cdot\frac{5}{2}=\frac{15}{8}=1\frac{7}{8}$ κιλά αλεύρι.
\sol{11} $\frac{2}{5}\cdot 1\frac{1}{2}=\frac{2}{5}\cdot\frac{3}{2}=\frac{3}{5}$ του λίτρου.
\sol{12} $\frac{3}{5}\cdot\frac{2}{3}=\frac{6}{15}=\frac{2}{5}$ του κήπου είναι ντομάτες. Σε $300$ τ.μ.: $\frac{2}{5}\cdot 300=120$ τ.μ.
\sol{13} $3\frac{1}{2}\cdot 2\frac{2}{5}=\frac{7}{2}\cdot\frac{12}{5}=\frac{84}{10}=\frac{42}{5}=8\frac{2}{5}$ τ.μ.
\sol{14} $4\frac{1}{2}\cdot\frac{2}{3}=\frac{9}{2}\cdot\frac{2}{3}=3$ χλμ.
\sol{15} Η έκπτωση είναι $\frac{3}{10}\cdot 80=24$ €. Πληρώνει $80-24=56$ €, που είναι και τα $\frac{7}{10}$ της τιμής: $\frac{7}{10}\cdot 80=56$.
\sol{16} Το μισό από το ένα τρίτο είναι $\frac{1}{2}\cdot\frac{1}{3}=\frac{1}{6}$, δηλαδή το ένα τρίτο χωρισμένο στη μέση. Είναι μικρότερο και από το $\frac{1}{3}$ και από το $\frac{1}{2}$. Ο πολλαπλασιασμός με κλάσμα μικρότερο του 1 δίνει λιγότερο από ό,τι είχες.
\sol{17} $12\cdot\frac{3}{4}=9$ (μικρότερο από 12) και $12\cdot\frac{5}{4}=15$ (μεγαλύτερο από 12). Όταν πολλαπλασιάζεις με κλάσμα μικρότερο του 1 ο αριθμός μικραίνει, και όταν πολλαπλασιάζεις με κλάσμα μεγαλύτερο του 1 μεγαλώνει.
\sol{18} Ο αντίστροφος του $\frac{3}{8}$ είναι ο $\frac{8}{3}$. Το κλάσμα που ψάχνουμε $x$ ικανοποιεί $x\cdot\frac{3}{8}=\frac{1}{2}$. Πολλαπλασιάζοντας και τα δύο μέλη με το $\frac{8}{3}$ βρίσκουμε $x=\frac{1}{2}\cdot\frac{8}{3}=\frac{4}{3}$. Έλεγχος: $\frac{4}{3}\cdot\frac{3}{8}=\frac{12}{24}=\frac{1}{2}$.
\sol{19} Οι αριθμητές και οι παρονομαστές απλοποιούνται διαδοχικά: $\frac{1\cdot 2\cdot 3\cdot 4\cdot 5}{2\cdot 3\cdot 4\cdot 5\cdot 6}=\frac{1}{6}$. Το γινόμενο ως το $\frac{9}{10}$ βγαίνει $\frac{1}{10}$.
\sol{20} Αν ο αριθμός $x$ ήταν αντίστροφος του 0, θα έπρεπε $0\cdot x=1$. Μα $0\cdot x=0$ για κάθε $x$, άρα η εξίσωση δεν έχει λύση. Το 0 δεν έχει αντίστροφο.
\sol{21} Με τις 4 στήλες και τις 3 σειρές το φύλλο έχει 12 ίσα μέρη. Η σκιασμένη περιοχή (3 στήλες) έχει 9 μέρη και από αυτά οι 2 σειρές είναι 6. Άρα τα σκούρα είναι $\frac{6}{12}=\frac{1}{2}$ του φύλλου. Και $\frac{2}{3}\cdot\frac{3}{4}=\frac{6}{12}=\frac{1}{2}$.
\sol{22} Κάθε «μισό του» διαιρεί με το 2: $80:2=40$, $20$, $10$, $5$. Άρα $\left(\frac{1}{2}\right)^4\cdot 80=\frac{80}{16}=5$.
